
\begin{proposition}[The irreducible local comparison and its rational and ruled consequences]\label{cdp:st:C.2.UN.7}
Let $f:Y\to B$ be a proper holomorphic map from a smooth
complex threefold to a smooth curve, with connected fibres
and smooth general fibre a complex two-torus. Let $mS$ be a full
fibre with irreducible reduced support, and let $F$ be a nearby
smooth fibre. Write $T_1$ for the local monodromy on
$H^1(F,\mathbb Q)$. Then
\[
 m\,e(S_{\rm reg})=\det(1-T_1^m)\ge0.
 \tag{C.UI.1}\label{cdp:eq:C.UI.1}
\]
For its normalization and minimal normalization resolution
$h=\eta\pi:Z\to T\to S$, the following statements hold.

\begin{enumerate}
\item If $Z$ is rational, then $e(S)\ge2$. More precisely,
      whenever $H^1(S,\mathbb Q)\ne0$, one has $e(S)=2$.
\item Suppose $Z$ is irrational ruled. A reduced horizontal
      conductor gives a smooth minimal elliptic ruled
      normalization, with disjoint smooth elliptic conductor
      curves and their finite etale double-cover maps to the
      conductor image. The support has ordinary double-curve
      normal crossings and Euler number zero.
\item Suppose the horizontal conductor is twice a noncontracted
      section. The normalization is again smooth minimal
      elliptic ruled. Its conductor is twice that section;
      its image is smooth elliptic, and the map from the
      section is etale of degree $1$, $2$, or $3$.
      The respective constant transverse types are the
      ordinary cusp, tacnode, and ordinary triple point.
      Every such support has Euler number zero.
\item Suppose the horizontal conductor is twice a section
      contracted by $\pi$. Then $e(S)\ge2$ and
      $\pi_1(S)=1$. In particular this case cannot occur
      under the full CDP hypotheses, since \ref{cdp:st:D.1.9} supplies
      a nonzero integral degree-one class on every fibre
      support there.
\end{enumerate}
The alternatives in the last three items exhaust the
irrational ruled case. No vanishing of $b_1(Y)$ or $b_2(Y)$
is used in the first four conclusions.
\end{proposition}
\begin{proof}
\leavevmode\par\medskip\noindent\textbf{The actual canonical and Chern-class comparisons.}\quad
Choose a coordinate $t$ on a small disc $\Delta$ about the
chosen value and put $U=f^{-1}(\Delta)$. The torsion-free
rank-one sheaf $f_*\omega_{U/\Delta}$ is a line bundle on
$\Delta$. Its rank is one by base change on the punctured
disc, where a two-torus has a one-dimensional space of
holomorphic top forms. A generator evaluates to a
holomorphic section of $\omega_{U/\Delta}$ with no zeros
on any smooth nearby fibre: on that fibre it is a nonzero
holomorphic two-form. Its divisor is therefore $kS$ for
some integer $k\ge0$. The coordinate $dt$ identifies the
relative canonical line with $K_U$. The actual fibre
equation and adjunction give
\[
 K_U\simeq\OO_U(kS),\qquad
 \OO_U(mS)\simeq\OO_U,\qquad
 \omega_S^{\otimes m}\simeq\OO_S.
 \tag{C.UI.2}\label{cdp:eq:C.UI.2}
\]
This restricts a section of the relative canonical line;
it does not pull a three-form back as a differential form
to a surface.

Here is the Euler calculation on this actual tube.
The map is flat: its local rings are torsion-free
over the regular one-dimensional base ring.
Proper flatness makes the holomorphic Euler
characteristic of its full fibres constant, so
$\chi(\OO_{mS})=\chi(\OO_F)=0$.
The Cartier filtration of $\OO_{mS}$ has successive
quotients $\OO_S(-jS)$, $0\le j<m$; by \eqref{cdp:eq:C.UI.2}
these are torsion line-bundle twists.
Each such twist has the same Euler characteristic
as $\OO_S$. One can check this without a compact
ambient manifold: use the normalization sequence
and the resolution Leray sequence to compute the
Euler characteristic on the compact smooth surface
$Z$, the conductor curve, and the finite-support
sheaf $R^1\pi_*\OO_Z$. Smooth-surface
Riemann--Roch leaves the first term unchanged
under a torsion twist; curve Riemann--Roch,
applied to the generic-rank filtration of the
conductor quotient, leaves the second unchanged;
the finite lengths in the third are unchanged.
These exact sequences prove the assertion on $S$.
Consequently, writing $c_i=c_i(TU)$, one has
\[
 \begin{split}
 0&=\chi(\OO_{mS})=m\chi(\OO_S),\\
 c_1(h^*TU)&=0\quad\hbox{in }H^2(Z,\mathbb Q),
 \qquad c_2(h^*TU)=0\quad\hbox{in }H^4(Z,\mathbb Q).
 \end{split}
 \tag{C.UI.3}\label{cdp:eq:C.UI.3}
\]
The first Chern-class assertion follows from \eqref{cdp:eq:C.UI.2}.
For the second we specify the local cycle comparison.
The inverse image of a radial segment $[0,t_0]$
is a compact triangulable real five-dimensional
space. Over the open segment it has the orientation
of the complex fibres followed by the radial
orientation. At a smooth point of $S$, the exact
local equation $t=u z^m$ gives $m$ radial sheets
with the same boundary orientation on $S$.
The remaining central singular locus has real
dimension at most two and contributes no
four-dimensional boundary. The resulting oriented
fundamental chain therefore has boundary
$[F]-m[S]$. This proves $[F]=m[S]$ in
$H_4(U,\mathbb Q)$, with its multiplicity retained.
The proper map $h$ has degree one onto $S$, hence
\[
 \int_Z h^*c_2=\langle c_2,[S]\rangle
       =\frac1m\int_F c_2(TU|_F)=0.
\]
The last equality follows from
$0\to T_F\to TU|_F\to\OO_F\to0$ and the
triviality of $T_F$. Top rational cohomology of
the connected closed oriented four-manifold $Z$
is one-dimensional, proving the second class
equality in \eqref{cdp:eq:C.UI.3}.

Finite duality and the minimal-resolution calculation
\ref{cdp:st:A.5.3} now give the actual effective integral divisor $D$
with
\[
                    K_Z(D)\simeq h^*\omega_S.
 \tag{C.UI.4}\label{cdp:eq:C.UI.4}
\]
Thus that line is torsion. Equations \eqref{cdp:eq:C.UI.2}--\eqref{cdp:eq:C.UI.4}
also give degree zero for both $K_Y$ and $\OO_Y(S)$
on every compact curve in this fibre. These are precisely
the ambient numerical inputs used below.

\par\medskip\noindent\textbf{The proper analytic Lefschetz formula.}\quad
Take an embedded resolution
$r:\widetilde U\to U$ of the pair, isomorphic over
$U\setminus\operatorname{Sing}S$, with a simple
normal-crossings central divisor
$\sum_i a_iE_i$. Its strict transform $E_0$ has
multiplicity $a_0=m$. Every exceptional multiplicity
satisfies
\[
                    a_i\ge2m\quad(i\ne0).
 \tag{C.UI.5}\label{cdp:eq:C.UI.5}
\]
Indeed a first centre outside the existing exceptional
locus lies in the singular locus of the reduced strict
transform, where its order is at least two. A centre
inside the existing exceptional locus already receives
the coefficient of an exceptional component, at least
$2m$. The coefficient of a new exceptional divisor is
the order of the current pullback of $f$ along its centre,
so the other contributions are nonnegative. Induction
proves \eqref{cdp:eq:C.UI.5}. The resolution can be chosen with all
centres above $\operatorname{Sing}S$. Its strict open
stratum is consequently
$E_0^\circ\simeq S_{\rm reg}$.

We specify the proper version of the trace formula used
here. A'Campo's original Theorem~1 is stated for a local
hypersurface germ. His \S2 begins with a proper analytic
map from a smooth analytic space to a disc, smooth off
the central simple-normal-crossings divisor, with no
Kahler assumption. The global boundary fibration and
its cohomology spectral sequence are constructed there;
the stalk-trace computation is on pp.~243--244
\cite[\S2 and \S3]{ACampo1975}. In the present notation
that construction gives
\[
 L(T^j):=\sum_q(-1)^q\operatorname{Tr}
       (T^j|H^q(F,\mathbb Q))
       =\sum_{a_i\mid j}a_i e(E_i^\circ).
 \tag{C.UI.6}\label{cdp:eq:C.UI.6}
\]
For clarity, the local and global maps producing this
identity are as follows. The real oriented blowup of
the divisor, followed by the angular map to $S^1$,
has a boundary fibre homeomorphic to $F$, as in
\ref{cdp:st:C.2.NC.1}. Above an $r$-branch stratum its fibre is
\[
 \{(u_1,\ldots,u_r)\in(S^1)^r:
                 u_1^{a_1}\cdots u_r^{a_r}=e^{i\theta}\}.
 \tag{C.UI.7}\label{cdp:eq:C.UI.7}
\]
An integral change of angular basis identifies this with
$\gcd(a_1,\ldots,a_r)$ copies of $(S^1)^{r-1}$.
Monodromy permutes these copies and acts by translations
on them. For $r\ge2$ the alternating trace is zero,
whether the permutation fixes a copy or not. For $r=1$
it is $a_1$ when $a_1\mid j$ and zero otherwise.
The constructible higher-direct-image sheaves of this
boundary projection carry those stalk actions. Their
cohomology spectral sequence computes $H^*(F)$ with
its actual monodromy. Triangulate the compact central
divisor compatibly with the strata and these local
systems. On every open cell the alternating trace of
compactly supported cochains is its signed cell
contribution times the stalk trace. Additivity and
the spectral sequence prove \eqref{cdp:eq:C.UI.6}. The proper
modification is an isomorphism over $\Delta^*$, so
these are the nearby fibre and monodromy of the
original map.

This description also proves quasi-unipotence without
an extra hypothesis. If $N$ is a common multiple of
the finitely many $a_i$, then $T^N$ is the identity
on every stalk cohomology group in \eqref{cdp:eq:C.UI.7}, since
the finite permutation disappears and translations
act trivially on cohomology. It is the identity on
the corresponding constructible cohomology sheaves,
hence on the second page of their cohomology spectral
sequence. The finite filtration on $H^q(F)$ has
identity action on its associated graded. Thus
$T^N$ is unipotent on $H^q(F)$.

Put $j=m$ in \eqref{cdp:eq:C.UI.6}. By \eqref{cdp:eq:C.UI.5} only $E_0$
contributes. The cohomology ring of the real
four-torus is the exterior algebra of its first
cohomology, and a monodromy diffeomorphism acts by
the exterior powers of $T_1$. Therefore
\[
 m e(S_{\rm reg})
   =\sum_{q=0}^4(-1)^q\operatorname{Tr}(\Lambda^qT_1^m)
   =\det(1-T_1^m).
 \tag{C.UI.8}\label{cdp:eq:C.UI.8}
\]
Every eigenvalue of $T_1^m$ is a root of unity by
the preceding quasi-unipotence. A real eigenvalue
is $1$ or $-1$ and contributes $0$ or $2$; each
nonreal conjugate pair contributes $|1-\lambda|^2$.
Their product is nonnegative. This proves \eqref{cdp:eq:C.UI.1}.

\par\medskip\noindent\textbf{The local degree-one class and the finite-character comparison.}\quad
We record the exact local replacement for the global
transgression used in \ref{cdp:st:C.2.MH.2}. Every nonzero class
$\gamma_S\in H^1(S,\mathbb Z)$ extends to a sufficiently
small tube and restricts nontrivially over $\mathbb Q$
to $F$. To see the extension, triangulate the analytic
pair $(Y,S)$ as in \ref{cdp:st:A.5.11} and take a regular
neighbourhood $N$ with a retraction $N\to S$.
Properness puts $f^{-1}(\Delta)$ inside $N$ after
shrinking $\Delta$. Pullback under that retraction
gives the specified class $\gamma_U$ on the tube,
whose restriction to $S$ is $\gamma_S$.

For injectivity upon restriction to $F$, the inclusion
$U^*=U\setminus S\hookrightarrow U$ surjects on
fundamental groups: loops can be perturbed off the
real codimension-two divisor, avoiding its higher
codimension singular strata. If a rational
degree-one class on $U$ restricts to zero on $F$,
the exact homotopy sequence of the punctured-disc
bundle makes its restriction to $U^*$ a scalar
multiple $a$ of the base-circle class. A small
meridian at a smooth point of $S$ is null-homologous
in $U$ and maps to $m$ times that base circle.
It follows that $ma=0$, hence $a=0$. Surjectivity
on fundamental groups now makes the original
class zero. In particular $\gamma_U|_F\ne0$.

Choose a prime $\ell$ for which both restrictions
to $S$ and $F$ stay nonzero modulo $\ell$.
Their values on one loop in each space exclude
only finitely many primes. The reduction of
$\gamma_U$ defines the actual cyclic etale cover
of $U$, with rational augmentation local system
$V$ of rank $d=\ell-1$. It is self-dual,
$H^0(S,V)=0$, and every nontrivial complex character
in it is nontrivial on $F$. The Koszul complex
of $\mathbb Z^4$ is then contracted by the inverse
of a nonzero character difference. Hence
\[
 H^*(F,V)=H^*(U^*,V)=0.
 \tag{C.UI.9}\label{cdp:eq:C.UI.9}
\]
The second equality is the bundle Leray comparison.

When $Z$ is rational or ruled over a curve of genus
$g$, all strata of an embedded resolution of the
central fibre are projective: the strict transform
is a blowup of $Z$, and every other surface arises
by blowups from a projective bundle over a compact
curve or from $\mathbb P^2$. \ref{cdp:st:C.2.MH.1} thus applies
to this actual local system. Its support comparison
makes $H^3(S,V)$ pure of weight three. The exact
normalization-conductor descent sequence has a
kernel which is a quotient of the degree-two
cohomology of a compact curve, hence pure of weight
two; strictness makes that kernel zero. The
negative-definite exceptional intersection matrices
then give the actual isomorphisms
\[
 H^3(S,V)\simeq H^3(Z,V),\qquad
 H^4(S,V)\simeq H^4(Z,V).
 \tag{C.UI.10}\label{cdp:eq:C.UI.10}
\]
These are precisely the maps (C.MH.10)--(C.MH.12);
none of their proofs uses ambient cohomology.
The acyclic punctured tube and local support duality
give, in the direct limit of shrinking tubes,
$H^j(S,V)\simeq H_{6-j}(S,V)$. Thus
\[
 d e(S)=2h^4(Z,V)-h^3(Z,V)
 =\begin{cases}
 2d,&Z\text{ rational},\\
 (2-2g)d,&Z\text{ ruled over genus }g.
 \end{cases}
 \tag{C.UI.11}\label{cdp:eq:C.UI.11}
\]
For the second line, the local system pulls back
from the ruling base. If its character is trivial
there, the two dimensions on the left are $d$
and $2gd$; otherwise they are $0$ and $(2g-2)d$.
This proves the precise local character version
of \ref{cdp:st:C.2.MH.2} and the additional assertion in item~1.

\par\medskip\noindent\textbf{Strict positivity for a rational normalization.}\quad
If $S$ is normal, \ref{cdp:st:C.2.UN.5} already makes it a
smooth torus or bielliptic surface, so its
resolution is neither rational nor irrational
ruled. Thus the remaining surface cases under
consideration are nonnormal.
The proof of \ref{cdp:st:A.5.16} uses the ambient assumption
only for $\omega_S$ torsion and $\chi(\OO_S)=0$,
now supplied by \eqref{cdp:eq:C.UI.2}--\eqref{cdp:eq:C.UI.3}. Consequently
its exact support and genus conclusions apply:
\[
 \begin{gathered}
 K_Z+D\simeq\OO_Z,\qquad D_{\rm red}\text{ connected},
 \qquad p_a(D_{\rm red})\le1,\\
 \Gamma_{\rm red}\text{ connected},\qquad
 p_a(\Gamma_{\rm red})\le2,\qquad
 e(\Gamma_{\rm red})\ge0,\\
 D_{\rm red}=h^{-1}(\Gamma_{\rm red}),\qquad
 S_{\rm reg}\simeq Z\setminus D_{\rm red}.
 \end{gathered}
 \tag{C.UI.12}\label{cdp:eq:C.UI.12}
\]
The final isomorphism retains all exceptional curves:
any exceptional fibre outside the conductor would
give an isolated normal hypersurface singularity,
excluded by the actual Milnor-sphere comparison
in \ref{cdp:st:A.5.17}. Thus
$e(S)=e(S_{\rm reg})+e(\Gamma_{\rm red})\ge0$.

Suppose it were zero. Both summands would be zero.
If $D_{\rm red}$ has an elliptic-normalized component,
the bound on its arithmetic genus makes that component
smooth and unique. Follow a sequence of rational
surface blowdowns to $\mathbb P^2$ or a Hirzebruch
surface. The divisor comparison at each step is
$D_{\rm new}=\sigma^*D_{\rm old}-E$; a centre of
multiplicity one increases the Euler number of the
complement by one, and a centre of multiplicity at
least two preserves it. A centre outside the
boundary is impossible by effectiveness. The
elliptic component survives these blowdowns.
On $\mathbb P^2$ the resulting anticanonical cubic
is a smooth elliptic cubic. On a Hirzebruch
surface the elliptic component is a smooth
bisection: a section is rational, and for a
coefficient-one bisection $H$ with vertical part
$V$ adjunction gives
$2p_a(H)-2+H\cdot V=0$, forcing $V=0$.
The minimal complement Euler numbers are therefore
$3$ and $4$, respectively. Their blowup comparison
contradicts $e(S_{\rm reg})=0$.

All normalized components of $D_{\rm red}$ are
therefore rational. Its first Betti number is at
most one, since the normalization sequence gives
$b_1(D_{\rm red})\le p_a(D_{\rm red})\le1$.
Every component of $\Gamma_{\rm red}$ is also
rational-normalized, being the finite image of a
strict conductor curve. Let their number be
$n\ge1$. For a connected compact curve with
rational normalizations the exact normalization
topology gives
$e(\Gamma_{\rm red})=1+n-b_1(\Gamma_{\rm red})$.
The supposed equality therefore gives
$b_1(\Gamma_{\rm red})=n+1\ge2$.

Here the required comparison with $S$ is an
actual pushout, not a presumed identification
of cohomology groups. By \eqref{cdp:eq:C.UI.12}, the natural
map from
$Z\cup_{D_{\rm red}}\Gamma_{\rm red}$ to $S$
is a continuous bijection from a compact space
to a Hausdorff space, hence a homeomorphism.
Triangulation makes $D_{\rm red}\hookrightarrow Z$
a cofibration, so its double mapping cylinder
gives the cohomology sequence of that pushout.
Connectedness and $H^1(Z,\mathbb Q)=0$ give
\[
 H^1(S,\mathbb Q)
  \simeq\ker\bigl[H^1(\Gamma_{\rm red},\mathbb Q)
                       \to H^1(D_{\rm red},\mathbb Q)\bigr].
 \tag{C.UI.13}\label{cdp:eq:C.UI.13}
\]
Its dimension is at least $n>0$. The proved
local finite-character calculation \eqref{cdp:eq:C.UI.11}
now gives $e(S)=2$, contradicting zero.
Thus $e(S)>0$, and \ref{cdp:st:C.2.UN.1} makes it an even
integer. This proves item~1.

\par\medskip\noindent\textbf{All noncontracted irrational ruled alternatives.}\quad
On a general ruling fibre $R$, \eqref{cdp:eq:C.UI.4} gives
$D\cdot R=2$. Thus its horizontal part is reduced
of total degree two, or is $2C$ for a smooth
section $C$. The proof of \ref{cdp:st:C.2.IRR.1} uses the
ambient condition only for the zero degrees
in (C.IR.2), (C.IR.7)--(C.IR.8) and for the
torsion-canonical, zero-$\chi$ genus bound \ref{cdp:st:A.5.14}.
All are furnished by \eqref{cdp:eq:C.UI.2}--\eqref{cdp:eq:C.UI.4}.
Its exact branch-reception, pinch-point and
blowdown comparisons therefore give item~2,
including its assertion that no horizontal
component is contracted.

Consider $D=2C+V$ with $C$ not contracted.
Every resolution-exceptional curve is vertical.
Indeed a positive-coefficient horizontal one
would have to be $C$. A coefficient-zero
exceptional curve $E$ satisfies
$K_Z\cdot E=D\cdot E=0$ by minimality and
adjunction. A horizontal exceptional curve over
a positive-genus base has $p_a(E)\ge1$ and
$E^2<0$, so $K_Z\cdot E>0$, a contradiction.
These are exactly the exceptional-curve
hypotheses of \ref{cdp:st:A.5.18}--\ref{cdp:st:A.5.20}.

Let $H$ be the reduced image of $C$ and
$d=\deg(C\to H^\nu)$. The branchwise conductor
formula is its intrinsic plane-branch conductor
plus the sum of its intersections with the
other branches. Its value on each horizontal
branch is two. If $d=1$, $H^\nu$ has the
positive genus of $C$, and no vertical rational
curve can dominate it. The sole branch with
conductor two is the ordinary cusp. If $d=2$,
the two horizontal branches must be smooth:
a singular branch would contribute at least two
intrinsically and at least one from the other
branch. Their allowed intersections give either
contact two between those branches, or simple
contact between them and with one additional
smooth branch. These are exactly a tacnode and
an ordinary triple point. For $d=3$ all three
branches must be smooth and pairwise transverse,
and there is no further branch. For $d\ge4$
the three other intersection contributions
already exceed two. This proves exhaustion
of the transverse cases used in the following
comparisons.

The local numerical hypotheses of \ref{cdp:st:A.5.14} are
exactly $\chi(\OO_S)=0$ and torsion $\omega_S$.
Its genus bound and delta-one complete-intersection
comparison consequently remain valid. The
conductor-cleared jets \ref{cdp:st:A.5.5} use degree zero for
$K_Y$ and $\OO_Y(S)$ on the image curve.
The cusp proof \ref{cdp:st:A.5.20} and the complete tacnode
proof \ref{cdp:st:A.5.18} therefore apply. The latter's
rational-image fundamental-group restriction
uses only the local quotient of
$\mathbb Z^4\rtimes\mathbb Z$, not an ambient
degree-one class. For the triple case, \ref{cdp:st:A.5.19}
excludes singular rational images with a
positive-genus covering component. \ref{cdp:st:C.2.D4.1}--2
exclude the connected degree-three and degree-two
covers of a smooth rational image. Their normal
cone is the actual cubic tensor in
$\operatorname{Sym}^3N_{H/Y}^*$; on
$H\simeq\mathbb P^1$ the needed identities
$\OO_Y(S)|_H=\OO_H$ and $\deg N_{H/Y}=-2$
follow from \eqref{cdp:eq:C.UI.2}. Their compact-core
transport is into the original nearby torus,
so no global Picard condition enters. The
independent ruled Chern comparison \ref{cdp:st:A.5.10} also
extends, since its additional input
$c_2(h^*TY)=0$ is \eqref{cdp:eq:C.UI.3}.

Consequently the image in each surviving case
is smooth elliptic and the transverse analytic
type is constant. No vertical rational curve
dominates that image, so all its generic
branches belong to $C$; the actual local
normalizations give etale degrees $1$, $2$,
or $3$, respectively. The normal line of $C$
has degree zero. The divisor comparison
(A.T.22), together with this square-zero
section and the absence of further conductor
along it, contracts no point of $C$ and makes
the minimal ruled boundary exactly $2C$.
A first reverse blowup away from $C$ would
create coefficient $-1$; a blowup on $C$
would create a conductor component meeting
it, contrary to the constant local type.
Thus the normalization is smooth minimal
elliptic ruled, exactly as proved in
\ref{cdp:st:A.5.18}--\ref{cdp:st:A.5.20}. The normalization Euler
comparison $e(S)=e(T)-e(C)+e(H)$ gives zero.
This proves item~3.

\par\medskip\noindent\textbf{A contracted section and its fundamental group.}\quad
Finally let $D=2C+V$ with $C$ contracted and
let $g\ge1$ be the ruling genus. The exact
boundary blowdown relation (C.CR.5), now
supplied by \eqref{cdp:eq:C.UI.4}, shows that $D_{\rm red}$
is the smooth section $C$ with trees of smooth
rational curves attached. All its crossings
are transverse nodes. Its inclusion in $Z$
induces an isomorphism on fundamental groups,
both identified with $\pi_1(C)$.
The proof of (C.CR.6) uses only this
connectedness, minimality and the isolated
singularity exclusion. It gives
\[
 D_{\rm red}=h^{-1}(\Gamma_{\rm red}),\qquad
 S_{\rm reg}\simeq Z\setminus D_{\rm red}.
 \tag{C.UI.14}\label{cdp:eq:C.UI.14}
\]
Every component of the connected curve
$\Gamma_{\rm red}$ is dominated by a vertical
rational curve and hence has normalization
$\mathbb P^1$. The same actual topological
pushout as above and van Kampen give
\[
 \pi_1(S)\simeq
 \pi_1(Z)*_{\pi_1(D_{\rm red})}\pi_1(\Gamma_{\rm red})
 \simeq\pi_1(\Gamma_{\rm red}).
 \tag{C.UI.15}\label{cdp:eq:C.UI.15}
\]
The second amalgamating map is trivial because
$C$ generates $\pi_1(D_{\rm red})$ and is
contracted. This group is free of rank
$\beta=b_1(\Gamma_{\rm red})$.

For its comparison with the nearby torus,
take the same regular neighbourhood and
small tube used above. A loop in $S$ lies
in the tube, and the neighbourhood retraction
is the identity on $S$. Thus the tube group
surjects onto $\pi_1(S)$. Perturbing loops
off $S$ makes it a quotient of
$\pi_1(U^*)=\mathbb Z^4\rtimes\mathbb Z$.
The latter is metabelian, whereas a free
group of rank at least two surjects onto
the nonmetabelian group $S_4$. Hence
$\beta\le1$. If $n\ge1$ is the number
of conductor-image components, its exact
normalization topology gives
\[
 e(\Gamma_{\rm red})=n+1-\beta\ge n\ge1,
 \qquad e(S)=e(S_{\rm reg})+e(\Gamma_{\rm red})>0.
 \tag{C.UI.16}\label{cdp:eq:C.UI.16}
\]
Here \eqref{cdp:eq:C.UI.1} proves the last inequality.
Parity makes $e(S)\ge2$. If $\beta=1$,
\eqref{cdp:eq:C.UI.15} supplies a nonzero integral
degree-one class on $S$; \eqref{cdp:eq:C.UI.11} then
gives $e(S)=2-2g\le0$, a contradiction.
Thus $\beta=0$, proving $\pi_1(S)=1$.
The transgression of \ref{cdp:st:D.1.9} excludes this
case under the full CDP hypotheses. This
proves item~4 and completes the proof.
\end{proof}

The assertion $e(S)=2$ in item~1 is deliberately tied
to the proved nonzero local-character comparison.
\ref{cdp:st:C.2.UN.8} constructs an irreducible rational-normalized
support with $\pi_1(S)=1$ and $e(S)=6$, so that
local character cannot be inserted without its
existence argument. The classification of all
irreducible zero-Euler supports is assembled in
\ref{cdp:st:C.2.UN.12} using this result, \ref{cdp:st:C.2.UN.5}, and the
class VII conductor comparisons. Reducible
unrestricted supports are not covered by this
irreducible statement.


